 \documentclass[book.tex]{subfiles}
\begin{document}

\section{Smooth scrolling on EGA}
\label{section:adaptive_tile_refresh}
Performing a full-screen redraw per frame would kill the CPU, as it would require updating every pixel on all four EGA planes. Achieving a 60Hz frame rate while refreshing the entire screen is impossible under these constraints. If we were to run the following code, which simply fills all memory banks, it would run at 5 frames per second.\\

\par
\begin{minipage}{\textwidth}
  \lstinputlisting[language=C]{code/write_ega.c}
  \end{minipage}
  \label{ega_write}

\bigskip
\par
\begin{fancyquotes}
Clearly, there is a whole world of awesome things there that we just couldn't do on the PC...
You can just redraw the whole screen, but then it turns out...\\
\par
Well, you're going five frames a second.\\

\par
\textbf{John Carmack\protect\footnotemark}
\end{fancyquotes}\\
\addtocounter{footnote}{-1}
\stepcounter{footnote}
\footnotetext{Interview with Lex Fridman in 2022, this quote is around 1h:41m.}

\par
So, how did they create a smooth scrolling game with these limitations? The solution was twofold:
\begin{itemize}
  \item By exploiting specific EGA hardware tricks.
  \item Updating only the portions of the screen that had actually changed between frames
\end{itemize}

\par
In the later episodes of Commander Keen, a different -- and significantly simpler -- approach was adopted for updating the display. Each of these solutions is described in detail in the following sections, starting with the EGA hardware.

\subsection{Moving one pixel at a time in EGA}
\label{section:EGA_virtual_screen}
The EGA card adds a powerful approach to linear addressing: the logical width of the virtual screen in VRAM does not need to match the physical screen display width. A programmer can define a virtual screen width of up to 4096 pixels and use the physical screen as a window onto any part of the virtual screen. What's more, the virtual screen's logical height can extend up to the total VRAM capacity. \\


\begin{figure}[H]
\centering
\includegraphics[width=1.0\textwidth]{imgs/drawings/virtual_screen.eps}
\caption{Virtual screen in VRAM.}
\label{fig:virtual_screen}
\end{figure}

The code below demonstrates how to set a custom logical width.\\

\begin{minipage}{\textwidth}
  \lstinputlisting[language=C]{code/SCREEN_WIDTH.ASM}
  \end{minipage}
  \label{ega_pel_pan}
  
\par
The displayed area of the virtual screen at any time is determined by setting the display memory address for fetching video data, configured via the CRTC Start Address register. The default address is \cw{A000:0000h}, though the offset can be adjusted to any other address.\\


\begin{minipage}{\textwidth}
  \lstinputlisting[language=C]{code/EGA_SET_ADDRESS.ASM}
  \end{minipage}
  \label{ega_set_address}
  \par


\par
To pan down a scan line, it only requires increasing the CRTC start address by the logical width.\\

\begin{figure}[H]
\centering
\includegraphics[width=1.0\textwidth]{imgs/drawings/virtual_screen_down.eps}
\caption{Move screen one pixel down.} 
\label{fig:crtc_down}
\end{figure}

\par
Horizontal panning is achieved by incrementing the start address by one byte. In EGA's planar graphics modes, the eight bits in each video RAM byte correspond to eight consecutive on-screen pixels, meaning the screen can move horizontally in steps of 8 pixels. This is rather coarse and not smooth horizontal scrolling. Fortunately, another EGA register addresses this limitation.\\

\begin{figure}[H]
\centering
\includegraphics[width=1.0\textwidth]{imgs/drawings/virtual_screen_right.eps}
\caption{Move screen eight pixels (1 byte) right.} 
\label{fig:crtc_right}
\end{figure}

\par
Smooth horizontal pixel scrolling is managed by the "Horizontal Pel Panning" register in the Attribute Controller (ATC), allowing fine adjustment in 1-pixel increments up to 7 pixels to the left. Programming the ATC requires attention: the ATC Index register only uses the lower five bits (bits 0-4) as its internal index. The next most significant bit, bit 5, controls the source of the video data sent to the monitor by the EGA card. When bit 5 is set to 1, the output of the color palette controls the displayed pixels; this is normal operation. When bit 5 is 0, video data doesn't come from the color palette, and the screen becomes a solid color. To maintain normal video operation, bit 5 must be set to 1 by writing \cw{20h} to the register.\\

\begin{minipage}{\textwidth}
  \lstinputlisting[language=C]{code/EGA_PELPAN.ASM}
  \end{minipage}
  \label{ega_pel_pan}


\trivia{"Pel" stands for "Picture element" and was introduced alongside pixel in the 1960s. While pixel became the standard term in graphics and modern computing, pel was frequently used in early image processing, computer technical references (notably IBM), and video coding.}\\

 
\par
By combining both EGA hardware tricks, smooth scrolling on EGA is achieved by adjusting the CRTC start address and fine-tuning the horizontal display using the horizontal pel panning register. The process consists of the following steps:
\begin{itemize}
  \item Calculate the panning offsets in pixels for both the x and y directions.
  \item Achieve smooth vertical scrolling by adding \textit{(logical width × y)} to the CRTC start address.
  \item For coarse horizontal scrolling in 8-pixel increments, increase the CRTC start address by \textit{x / 8} bytes.
  \item Finally, apply fine horizontal pixel alignment by writing \textit{(x \% 8)} to the horizontal pel panning register.
\end{itemize}

\begin{figure}[H]
\centering
\includegraphics[width=0.7\textwidth]{imgs/drawings/Tile_Refresh.eps}
\caption{Smooth scrolling in EGA.}
\label{fig:tile_refresh}
\end{figure}

\par
The engine is now capable of moving the entire screen one pixel at a time in all directions, just by tweaking the EGA registers, enabling smooth scrolling across the display.\\




\subsection{Adaptive Tile Refresh}
So far, we have built a virtual screen in VRAM allowing smooth one pixel moves in both axes. However, once the virtual screen edge is reached, a full-screen redraw would be too slow, causing a drop to 5 fps.\\

\par
A naive approach for screen refresh is to load the entire level into VRAM, setting the logical width to match the level width. By adjusting the CRTC Start Address and Horizontal Pel Panning registers, smooth scrolling can be achieved. unfortunately, loading the entire level requires far more VRAM than is available. The first level, \textit{Horse Radish Hill}, is 136 tiles wide and 37 tiles high. Each tile is 128 bytes,  which means a total of 136$\times$37$\times$128 bytes = 644KB is required to store the entire level in VRAM. Since we only have 256KiB available, this is not a feasible option.\\

\textbf{\underline{Trivia :}} At the time of writing this book, most video cards contain more than 1GB VRAM,  sufficient to store all Commander Keen levels at once in VRAM. \\

\par
To solve the refresh issue, John Carmack invented "Adaptive Tile Refresh" (ATR). The core idea is to refresh only those areas of the screen that need to change\footnote{See https://retrocomputing.stackexchange.com/questions/22175/what-is-adaptive-tile-refresh-in-the-context-of-commander-keen}. Let's look at \textit{Commander Keen 1: Marooned on Mars} in Figure \ref{fig:keen_difference} on the next page. This is the first level of Marooned, immediately to the right of the crashed Bean-with-Bacon Megarocket. The first figure is the start of the level, the second figure is after Keen has scrolled one tile (16 pixels) to the right through the world. They look almost identical to the naked eye, don't they? \\

\par
Now, if we perform a difference on both images, you can see which tiles need to be changed upon screen refresh. The trick behind the scrolling is to only redraw tiles that actually changed after panning 16 pixels (one tile). For matching tiles there is nothing to do and they are skipped entirely. In total, only 69 tiles out of the 260 tiles need to be refreshed, which is 27\% of the screen!\\

\par
\begin{fancyquotes}
The trick behind the scrolling in the first Commander Keen games was to only redraw tiles that actually changed after panning 16 pixels, since most maps had large swathes of constant background.\\
\par
Text consoles have similar statistics with spaces, so the same trick is a big win.\\
\par
\textbf{John Carmack\protect\footnotemark}
\end{fancyquotes}\\
\addtocounter{footnote}{-1}
\stepcounter{footnote}
\footnotetext{John Carmack on X, July 11, 2020.}

\par
This is also the part where the game designers Tom Hall and John Romero had to help the engine. Smooth scrolling is inversely proportional to the number of tiles to redraw. A checkerboard tile pattern basically means a full-screen redraw and would kill the CPU. So, to avoid costly "jolts", the game designers built tile maps with a lot of repeating tiles. \\ 

\begin{figure}[H] 
  \centering 
  \includegraphics[width=.66\textwidth]{screenshots_300dpi/game/keen-ATR-diff.png}
  \caption{Start of the world, moved one tile to the right and difference.}
  \label{fig:keen_difference}
\end{figure}
\pagebreak


\par
Let's have a closer look at the EGA VRAM setup. The video memory is organized into three virtual screens:
\begin{itemize}
\item Page 0 and 1, which are used to switch between buffer and visible screen. The idea between two pages (double buffer) is that the code can draw in the second buffer while the first buffer is being shown on screen, which is then switched out during screen refresh. This ensures that no frame is ever displayed mid-drawing, which yields smooth, flicker-free animation.
\item A master page containing a static page, which is copied to the buffer screen when performing the screen refresh.
\end{itemize}


\begin{figure}[H]
\centering
\includegraphics[width=\textwidth]{imgs/drawings/ega_ram_architecture.eps}
\caption{Virtual screen layout on EGA card.}
\label{fig:ega_ram_arch}
\end{figure}

\par
The EGA screen in mode \cw{0Dh} has a resolution of 320x200 pixels, or 40x200 bytes. Let's extend the height by 8 bytes to have a height of 208 pixels, so the screen fits nicely in 20x13 tiles. By making the virtual screen one tile higher and one wider on each side of the screen, the engine can scroll up to 16 pixels to any direction of the screen without any tile refresh, by simply adjusting the CRTC Start Address and Pel Pan registers.\\

\par
Each virtual screen has a size of 44$\times$240$\times$4 = 42,420 bytes. So, within a 256KiB EGA card there is enough VRAM available to keep all three virtual screens in memory. The page that is displayed on screen is selected by setting the CRTC Start Address register at which to begin fetching video data.\\

\begin{figure}[H]
\centering
\includegraphics[width=1.0\textwidth]{imgs/drawings/ATR_virtual_screen.eps}
\label{fig:atr_virtual screen}
\end{figure}


\par
For both the buffer and visible screen a tile array is created to maintain which tiles are updated since last refresh.\\


\par
\begin{minipage}{\textwidth}
  \lstinputlisting[language=C]{code/tile_index_array.c}
  \end{minipage}
  \label{ega_refresh}
  \par

The steps to refresh the screen are as follows (ignore sprites for now):
\begin{enumerate}
\item Check if the player has moved one tile in any direction.
\item The ATR function scans through every tile in the current display page and compares them against the corresponding tile from the level map one by one.
\item If the stored tile index number doesn't match the one in the level map, the corresponding tile is copied to the master page, and marked in both tile update arrays. 
\item Refresh the buffer page by copying marked tiles from the master page.
\item Switch the view and buffer pages by adjusting the CRTC Start Address and Pel Panning registers.
\end{enumerate}

\par
In the next three snapshots, we take you step-by-step through each of the stages. The player has reached the edge of the virtual screen and moves further to the right.\\ 

 
\begin{figure}[H]
\centering
 \includegraphics[width=0.84\textwidth]{screenshots_300dpi/game/Keen_ATR_1-3_a.png}
 \caption{Start: Reach the end of the virtual screen.}
 \label{fig:kc1_3_start}
\end{figure}

\pagebreak
The engine keeps track of which tile numbers are part of the virtual screen. Since the engine only refreshes the screen at tile size granularity, it can determine extremely fast what has changed on the screen by comparing tile numbers. If the tile number has changed, the tile is updated by copying tiles from RAM into the VRAM master page. The changed tiles are marked with a '1' in both tile arrays, meaning it needs to be updated upon next refresh.

\begin{figure}[H]
\centering
 \fullimage{/game/Keen_ATR_1-3_step1.png}
 \caption{Update tiles in master page and update both tile update arrays.}
 \label{fig:kc1_3_step1}
\end{figure}

\par
The next step is to scan all tiles in buffer tile update array (array 1) and for each tile marked '1', copy the corresponding tile from master to virtual page 1 page.\\

\begin{figure}[H]
\centering
 \fullimage{/game/Keen_ATR_1-3_step2.png}
 \caption{Copy changed tiles from virtual master page to page 1.}
 \label{fig:kc1_3_step2}
\end{figure}


\pagebreak
In the final step, point the visible screen to virtual page 1 by updating the CRTC start address. Finally, tile update array 1 is cleared to '0'. Now the first step is repeated, but this time virtual page 0 acts as the buffer screen. Note that after swapping, tile update array 0 keeps marked tiles from last update. This makes sense, as the current buffer page is not yet refreshed since it was displayed in the previous refresh cycle. 

\begin{figure}[H]
\centering
 \fullimage{/game/Keen_ATR_1-3_step3.png}
 \caption{Update CRTC Start Address to virtual page 1 and empty tile update array 1.}
 \label{fig:kc1_3_step2}
\end{figure}

\pagebreak

Now the buffer is refreshed and the CRTC address is updated, the final step is fine adjustment using the Pel Panning register.\\
\begin{figure}[H]
\centering
 \includegraphics[width=0.84\textwidth]{screenshots_300dpi/game/Keen_ATR_1-3_b.png}
 \caption{Screen is refreshed and scrolled to the right.}
 \label{fig:kc1_3_start}
\end{figure}



\subsection{Virtual Screen Tile Refresh}
\label{section:wrap_ega_memory}
\label{section:optimize_tile}
John Carmack explored what would happen if you push the virtual screen over the 64KiB border (address \cw{FFFFh}) in video memory. It turned out that the EGA continues the virtual screen at \cw{0000h}. This means you could wrap the virtual screen around  the EGA memory and only need to add a stroke of tiles on one of the edges when Commander Keen moves more than 16 pixels.\\
\par
\begin{figure}[H]
  \centering
  \includegraphics[width=0.95\textwidth]{imgs/drawings/ega_wrapping.eps}
  \label{fig:ega_wrapping}
  \caption{EGA memory wrap-around.}
\end{figure}
 
\par
\begin{fancyquotes}
I finally asked what actually happens if you just go off the edge [OF THE VRAM]?\\

If you take your [CRTC] start and you say OK, I can move over and I get to what should be the bottom of the memory window. [...] What happens if I start at 0xFFFE at the very end of the 64k block? It turns out it just wraps back around to the top of the block.\\

I'm like oh well this makes everything easy. You can just scroll the screen everywhere and all you have to draw is just one new line of tiles.\\

It just works. We no longer had the problem of having fields of similar colors. It doesn't matter what you're doing, you could be having a completely unique world and you're just drawing the new strip.\\
\par
\textbf{John Carmack\protect\footnotemark}
\end{fancyquotes}\\
\addtocounter{footnote}{-1}
\stepcounter{footnote}
\footnotetext{An explanation further elaborated during the same interview with Lex Fridman in 2022.}




\subsection{Virtual Screen Tile Refresh in Keen Dreams} 
\label{section:scroll_refresh_dreams}
The EGA memory wrapping results in an improved, more simplified screen refresh algorithm. This algorithm was called "Virtual Screen Tile Refresh".\\

\begin{fancyquotes}
The second Keen trilogy used a better trick -- just keep panning and redrawing the leading edge, letting the screen wrap around at the 64k aperture edge.\\
\par
\textbf{John Carmack\protect\footnotemark}
\end{fancyquotes}\\
\addtocounter{footnote}{-1}
\stepcounter{footnote}
\footnotetext{John Carmack on Twitter.}

\par
First, a virtual screen and tile update array is defined by making the display one tile taller and wider, creating what is known as the viewport. This allows the engine to scroll the display up to 16 pixels to the right and bottom without refreshing tiles, by simply adjusting the CRTC Start Address and Pel Pan register. \\




\begin{figure}[H]
\centering
\includegraphics[width=\textwidth]{imgs/drawings/buffer_tile_layout.eps}
\caption{Tile update array layout.}
\label{fig:screen_setup}
\end{figure}

\pagebreak
An additional column is then added to the viewport to support horizontal scrolling. Finally, the tile update array is further expanded to allow the entire viewport to move up to two tiles in any direction.\\

\par
The full refresh cycle, including sprite updates, proceeds as follows:
\begin{enumerate}
  \item Verify if the player has moved one tile in any direction.
  \item Update both the tile update array and VRAM pointers, copy the new column or row of tiles to the master page and mark the tiles in both arrays.
  \item Refresh the buffer page by scanning the tile index array. If a tile is marked, copy it from the master page to the buffer page.
  \item Iterate through the sprite removal list, copying the corresponding image block from the master page to the buffer page to clear the sprite.
  \item Iterate through the sprite list, copying each sprite image block to the buffer page.
  \item Switch the view and buffer pages by adjusting the CRTC Start Address and Pel Panning registers.
\end{enumerate}


\vspace{0.5cm}
\par
\begin{minipage}{\textwidth}
  \lstinputlisting[language=C]{code/refresh.c}
\end{minipage}



\pagebreak
Let's go over the refresh cycle step by step. At the start, all three virtual pages display the same viewport, with virtual page 0 currently shown on-screen. Both tile update arrays are empty, meaning no tile updates are required in the next refresh cycle. The player has reached the edge of the virtual screen and moves further to the left. \\

\begin{figure}[H]
\centering
 \scaledimage{1.0}{/game/Keen_ATR_4-6_a.png}
 \caption{VRAM and tile update arrays before scrolling to the left.}
 \label{fig:kc4_6_step0}
\end{figure}

\par
Both VRAM \cw{screenstart[]} pointer and tile update array pointer shift left to introduce a new column of tiles. The new column of tiles is copied from RAM to the master page, while the leftmost column in both tile update arrays is marked with '1' to ensure it updates in the next refresh cycle. Animated tiles are then copied to the master page, and the corresponding tiles in the array are also marked with '1'.\\


\vspace{-5pt}
\begin{figure}[H]
\centering
 \scaledimage{1.0}{/game/Keen_ATR_4-6_b.png}
 \caption{Update Virtual master with new left column and animated tiles.}
 \label{fig:kc4_6_step1}
\end{figure}

\pagebreak
Next, the engine scans all '1's and copies the corresponding tiles from the master to the buffer page. \\

\vspace{-5pt}
\begin{figure}[H]
\centering
 \scaledimage{1.0}{/game/Keen_ATR_4-6_c.png}
 \caption{Copy changed tiles from master to buffer page.}
 \label{fig:kc4_6_step2}
\end{figure}


\par
The final step is to erase all sprites listed in the sprite removal list from the buffer page and copy new sprites from RAM to the buffer page. Tiles overlapping with erased blocks are marked with '2' in the tile update array, while tiles containing new sprites are marked with '3'. Note that the '2' marker is not used anywhere in the engine; it was most likely part of the original ATR algorithm.\\

\begin{figure}[H]
\centering
 \scaledimage{1.0}{/game/Keen_ATR_4-6_d.png}
 \caption{Removing and updating sprites to the buffer screen.}
 \label{fig:kc4_6_step4}
\end{figure}

\par
For each sprite that needs to be removed from the buffer screen, its location and size are stored in the sprite removal list, allowing the sprite to be removed by copying the corresponding section from the master screen to the virtual page.\\
\par
\begin{minipage}{\textwidth}
  \lstinputlisting[language=C]{code/erase_blocks_struct.c}
  \end{minipage}
  \label{block_removal}
\\


\par
Notice that the removal block is perfectly aligned on an even memory address by using the \cw{\& \textasciitilde{}1} bitwise operator to clear the least significant bit, ensuring 16-bit transfers per CPU cycle.\\

\par
\begin{minipage}{\textwidth}
  \lstinputlisting[language=C]{code/erase_blocks.c}
  \end{minipage}
  \label{block_removal}


  


\pagebreak
Once the buffer refresh is completed, the engine needs only to update the CRTC start address to virtual page 1. At this point, the visible tile update array is emptied, and the pointer resets to its starting location. The new buffer array (virtual page 0) retains the marked tiles since this screen has not yet been refreshed. Finally, the refresh cycle restarts, with virtual page 0 now functioning as the buffer. 

\begin{figure}[H]
\centering
 \scaledimage{0.99}{/game/Keen_ATR_4-6_e.png}
 \caption{Swap buffer and visible screen by updating CRTC start address.}
 \label{fig:kc4_6_step4}
\end{figure}

\pagebreak
The final step in scrolling is to adjust the horizontal fine-pixel alignment by setting the Pel Pan register. \textit{Et voil\`a}, the screen scrolls smoothly to the left.\\

\begin{figure}[H]
\centering
\begin{subfigure}{.5\textwidth}
\centering
  \includegraphics[width=.95\linewidth]{screenshots_300dpi/game/Scroll_KC4_6_1.png}
\end{subfigure}%
\begin{subfigure}{.5\textwidth}
\centering
  \includegraphics[width=.95\linewidth]{screenshots_300dpi/game/Scroll_KC4_6_2.png}
\end{subfigure}
\caption{Left screen is before scrolling, right screen after scrolling.}
\label{fig:kc4_6_scroll}
\end{figure}

\par
If the player continues to move to the left, the tile update array 0 pointer lowers a second time, underscoring why the viewport needs to float up to two tiles in all directions.\\

\par
This scrolling process is significantly more efficient than the original adaptive tile refresh engine, requiring only 8\% of the tiles to be refreshed. In addition, level design is no longer constrained by the need to use large repeating patterns.\\

\par
This is illustrated on the next page with a comparison between a level from Commander Keen II - \textit{The Earth Explodes} (top) and one from Keen IV - \textit{Secret of the Oracle} (bottom). Notice how in Commander Keen IV almost nothing repeats except for the minimal path and underground sections.\\

\begin{figure}[H]
\centering
\begin{subfigure}[b]{\textwidth}
  \fbox{\includegraphics[width=1.0\textwidth]{screenshots_300dpi/keen2.png}}
\end{subfigure}
\par\bigskip\bigskip % force a bit of vertical whitespace
\begin{subfigure}[b]{\textwidth}
  \fbox{\includegraphics[width=1.0\textwidth]{screenshots_300dpi/keen4.png}}
\end{subfigure}
\end{figure}




\pagebreak
\subsection{Crippled Super VGA compatibility}
\par
The introduction of Super VGA cards created a compatibility issue, as these cards typically contained more than 256KiB of VRAM\footnote{In 1989 the VESA consortium standardized an API to use Super VGA modes in a generic way. One of the first modes was 640x480 at 256 colors requiring more than 256KiB VRAM, which from a hardware constraint resulted in 512KiB.}. This resulted in crippled backwards compatibility and the wrapping around \cw{FFFFh} did not work anymore on these cards. \\

\par
\begin{fancyquotes}
I was in a tough position. Do I have to track every single one of these [SUPER VGA CARDS] and it was a madhouse back then with 20 different video card vendors with all slightly different implementations of their non-standard functionality. Either I needed to natively program all of the cards or I kind of punt. I took the easy solution of when you finally did run to the edge of the screen I accepted a hitch and just copied the whole screen up.\\
\par
\textbf{John Carmack\protect\footnotemark}
\end{fancyquotes}\\
\addtocounter{footnote}{-1}
\stepcounter{footnote}
\footnotetext{And again the same interview with Lex Fridman in 2022.}


\par
A simple solution can be used to resolve this issue. As illustrated in Figure \ref{fig:ega_ram_arch} on page \pageref{fig:ega_ram_arch}, the space between \cw{B400h} and \cw{FFFFh} is unused and provides enough room for another virtual screen. Instead of relying on memory wrapping, the corresponding screen page is copied to the opposite end of the VRAM buffer, as illustrated in Figure \ref{fig:page_wrapping}. This ensures that the display remains contiguous.\\


\par 
\begin{figure}[H]
\centering 
\begin{subfigure}{.5\textwidth}
  \centering
  \includegraphics[width=.8\textwidth]{imgs/drawings/Page_down_switch.eps}
  \caption*{Copy Page 0 on top of Master Page.}
  \label{fig:page0_down}
\end{subfigure}% 
\begin{subfigure}{.5\textwidth}
  \centering
  \includegraphics[width=.8\textwidth]{imgs/drawings/Page_up_switch.eps}
    \caption*{Copy Page 0 below Page 1. }
  \label{fig:page0_up}
\end{subfigure}
\caption{Move screen to opposite end of VRAM buffer}
\label{fig:page_wrapping}
\end{figure}



\par
\begin{minipage}{\textwidth}
  \lstinputlisting[language=C]{code/wrap_page_buffer.c}
  \end{minipage}
  \label{ega_refresh}
  \par
As explained in Section "\nameref{section:ega_memmap}" on page \pageref{section:ega_memmap}, each pixel is encoded by four bits, which are distributed across the four EGA banks. Copying the entire page requires reading all data from one EGA bank into RAM, writing it back to VRAM at its new location, and repeating this process for each of the four EGA banks. So, how can we copy four VRAM planes fast enough without introducing a noticeable performance hit? \\  

\par
A closer look at the EGA card reveals that a latch is placed in front of the ALU, which can be repurposed for this task. The ALU for each bank can be configured to use only the latch value when writing to VRAM, ignoring the CPU-provided value\footnote{The mask trick is discussed in \textit{Game Engine Black Book: Wolfenstein 3D} for the VGA card. The trick works the same way on the EGA card.}. By enabling the mask to write to all four memory banks, the values stored in the latches are written simultaneously to VRAM. With this setup, a single read operation \circled{1} loads all four latches at once, and \circled{2} four bytes—one per bank—are written with a single write operation. This makes it possible to copy the entire buffer fast enough without causing any noticeable performance impact.


\par
 \begin{figure}[H]
\centering
 \includegraphics[width=\textwidth]{imgs/drawings/latches.eps}
 \caption{Read operation (in red) stored in latches which are later used to write back into VRAM, ignoring the RAM value (in blue).}
 \label{fig:latch_memory}
 \end{figure}
 
\par
The setup requires to configure the graphics controller to only read the content of the latches and the map mask register to enable writing to all 4 planes at once.\\ 

\par
\begin{minipage}{\textwidth}
  \lstinputlisting[language=C]{code/EGA_LATCH_COPY.ASM}
  \end{minipage}
  \label{ega_latch_copy}

\subsection{Screen refresh} 
Flipping between the pages is as simple as setting the CRTC start address registers to page 0 or page 1 starting point. However, there is one issue to solve. If you were to run it, every once in a while the expected screen shown below...

\setlength{\fboxsep}{0pt}
 \begin{figure}[H]
\centering
 \fbox{\includegraphics[width=0.95\textwidth]{screenshots_300dpi/game/full_screen.png}}
 \end{figure}
 \vspace{-8pt}
...would instead appear distorted, showing both misalignment and parts of two pages:
 \begin{figure}[H]
\centering
 \fbox{\includegraphics[width=0.95\textwidth]{screenshots_300dpi/game/crtc_scanstart_problem.png}}
 \end{figure}
\par
This problem results from the timing between updating the CRTC start address and the screen refresh. The start address is latched by the EGA's internal circuitry exactly once per frame, usually at the beginning of the vertical retrace. Although the CRTC start address is a 16-bit value, the \cw{out} instruction can only write 8 bits at a time.\\

\par
This issue can be illustrated with the following setup: the current CRTC start address (Page 0) points to \cw{0000h}, while the buffer (Page 1) points to \cw{3C00h}. After moving one tile to the left, Page 0 now points to \cw{FFFEh} in VRAM, and Page 1 points to \cw{3BFEh}. Since Page 1 is the updated buffer, it will be displayed in the next refresh cycle. However, due to poor timing in updating the vertical retrace and start address, the CRTC only picks up the first byte of the address, \cw{3Bh}, setting the start address to \cw{3B00h} instead of \cw{3BFEh}.\\

\bigskip
\begin{minipage}{\textwidth}
\lstinputlisting[language={[x86masm]Assembler}]{code/pageflip_error.c}
\end{minipage}

\par
One approach to this problem is to update the start address when the vertical retrace signal is detected via the Input Status 1 Register, specifically bit 3 at port \cw{3DAh}. Unfortunately, by the time the program observes this status, the start address for the next frame has already been latched, as this occurs immediately at the beginning of the vertical retrace pulse.\\

\par
An alternative approach is to update the CRTC start address at a point sufficiently distant from the start of the vertical retrace. The Display Enable Status signal at bit 1 from the same Input Status Register can be used for this purpose; a value of 1 indicates that the display is currently within a horizontal or vertical retrace\footnote{Documentation is a bit unclear here. The IBM technical documentation for VGA explains retrace takes place when bit 0 of the Input Status Register 1 is set to high. The IBM technical EGA documentation explains the opposite, saying when bit 0 is set low a retrace is taking place. For now, we assume source code and VGA documentation is correct, retrace takes place on a '1'.}. What we want to detect is the completion of a horizontal or vertical retrace and the beginning of a new scan line, which is far enough from the vertical retrace to ensure that the new start address is latched during the next vertical sync.\\

\par
To guarantee that the start address is updated at the very beginning of a new scan line, it first waits for the current scan line to complete. It then waits for a full retrace, verifying that the Display Enable Status returns to 0. At this point, the CRTC start address is updated. Once the CRTC start address has been updated, the engine waits for a vertical retrace to set the Pel panning register.\\

 
\begin{figure}[H]
  \centering
  \includegraphics[width=\textwidth]{imgs/drawings/update_start_address.eps}
  \label{fig:update_start_address}
  \caption{Update CRTC start address at beginning of new scan line.}
\end{figure}


\trivia{Regardless of CPU speed, the game's speed is limited by the monitor's refresh rate, which is 60Hz for EGA and 70Hz for VGA.}\\


Wolfenstein 3D is using an even faster solution. Here, each virtual screen has a fixed starting address and only differ by their high byte value: page 0 is at \cw{0000h}, page 1 at \cw{4100h} and the master page at \cw{8200h}. Moving from any page to another requires updating only the high 8 bits, ensuring a smooth screen refresh\footnote{See \textit{Game Engine Black Book: Wolfenstein 3D} section "Solving the VGA Problem".}. This solution is not possible for Commander Keen as the virtual screen moves around in VRAM. 

\par
\begin{minipage}{\textwidth}
\lstinputlisting[language={[x86masm]Assembler}]{code/screen_refresh.c}
\end{minipage}
\par


\section{Actors and AI}
\label{section:actors_and_ai}
All objects in the game, such as Commander Keen, enemies, bonus points, doors, and even the scoreboard, are called "actors". Each actor is controlled via a state machine, enabling them to think and take actions such as walking, throwing objects, jumping over a gap, or playing sound. To model their behavior, each enemy has an associated state, which can include:
\begin{itemize}
  \item Chasing Keen
  \item Attacking or smashing Keen
  \item Shooting projectiles
  \item Climbing and sliding on poles
  \item Turning into a flower
  \item Special Boss (Boobus)
\end{itemize}

\subsection{State Machine}
A state machine is a design pattern that manages an object's behavior by breaking it into distinct states (e.g., idle, running, jumping) and defining rules for transitioning between them.
In Commander Keen, each state has associated think, reaction, and contact method pointers. Additionally, there is a \cw{tictime} value and \cw{*nextstate} pointer, which indicate when the actor should transition to another state after a specific number of tics have passed in the current state.\\

\par
\begin{minipage}{\textwidth}
  \lstinputlisting[language=C]{code/statetype.c}
\end{minipage}
\label{state_type}
\par

Each actor has a defined state chain, for example the Pea Pod. \\
\par
\begin{minipage}{\textwidth}
\lstinputlisting[language=C,style=mystyle,basicstyle=\small]{code/s_grdchase.c}
\end{minipage}

\par
Different types of enemies have their own state machines, often sharing reaction functions (e.g., \cw{WalkReact} and \cw{ProjectileReact}) but usually having their own unique "thinking" states.\\

\begin{figure}[H]
  \centering
  \includegraphics[width=1.0\textwidth]{imgs/drawings/state_machine.eps}
  \caption{State machine for Pea pod and Pea brain.}
  \label{fig:state_machine}  
\end{figure}

\par
The \cw{*react} function determines how an enemy reacts to the environment, such as turning around when hitting a wall or reaching the edge of a platform. The \cw{*think} function defines how an enemy behaves when Commander Keen is nearby (e.g., attacking or firing a projectile) or when it reaches an edge (e.g. jumping). In some cases it introduces randomness, like when a pea pod might decide to spit a pea brain.

\par
\begin{minipage}{\textwidth}
  \lstinputlisting[language=C]{code/state_machine_think.c}
\end{minipage}

\par
The \cw{*contact} function checks if an object has come into contact with another object and defines the resulting interaction, such as Commander Keen losing a life.



\subsection{Clipping}
\label{section:clipping}
Whether an actor can move through or fall through a tile is determined by the tile property \cw{tinf[]}. Each foreground tile includes four directional parameters: north, south, east, and west. If, for example, the east parameter has a value greater than 0, the tile has a solid east wall, and the actor cannot enter it from the east side. When an actor moves to the left and encounters a solid tile from the east, the engine adjusts the actor's movement to prevent it from crossing the tile.\\


\begin{figure}[H]
  \centering
  \includegraphics[width=\textwidth]{imgs/drawings/clipping_east.eps}
  \caption{Clipping to east wall when actor moves left.}
  \label{fig:clipping_east}  
\end{figure}

\par
\begin{minipage}{\textwidth}
  \lstinputlisting[language=C]{code/clip_east_wall.c}
\end{minipage}
\label{wallclip_array}
\par
For clipping along the top and bottom of tiles, the engine also accounts for standing on slopes. When the actor is clipped to the top or bottom of a slope tile, an offset is applied to move the actor up or down along the slope. \\

\begin{figure}[H]
  \centering
  \includegraphics[width=\textwidth]{imgs/drawings/clipping_north.eps}
  \label{fig:clipping_north}
  \caption{Clipping north wall with slope.}
\end{figure}

This offset is defined by the lookup table \cw{wallclip[][]}, which uses the actor's midpoint and the wall type to determine the slope type.

\par
\begin{figure}[H]
  \centering
  \includegraphics[width=\textwidth]{imgs/drawings/walltype.eps}
  \caption{\cw{wallclip[8][16]}: Fall through, solid and 6 slope types.}
  \label{fig:walltype}
\end{figure}



\par
\begin{minipage}{\textwidth}
  \lstinputlisting[language=C]{code/clip_north_wall.c}
\end{minipage}


\section{Rendering the layers}
\label{section:DrawingLayer}
Each tile on the screen can contain up to three layers: the background tile layer, the foreground tile layer and a sprite layer. Rendering the final result on screen requires multiple redraws of the screen:
\begin{enumerate}
  \item Draw the background tile or a combined background and foreground tile.
  \item Draw the sprites.
  \item Re-draw the foreground tile if a sprite should not appear on top.
\end{enumerate}

\par
It must run quickly and therefore most of the code is written in assembly.

\subsection{Draw background and foreground tiles}
Drawing the background layer is straightforward: it only requires copying 128 bytes to VRAM. Drawing a foreground tile on top of the background tile requires an additional mask, which defines which pixels are overwritten by the foreground tile.\\

\begin{figure}[H] 
\centering
 \begin{subfigure}{0.32\textwidth}
 	\centering
 	\includegraphics[width=0.8\textwidth]{screenshots_300dpi/game/mask_tiles_1.png}
 	\caption{Background tile.}
 \end{subfigure}
 \begin{subfigure}{0.32\textwidth}
 	\centering
 	\includegraphics[width=0.8\textwidth]{screenshots_300dpi/game/mask_tiles_2.png}
 	\caption{\cw{AND} foreground mask.}
 \end{subfigure}
 \begin{subfigure}{0.32\textwidth}
 	\centering
 	\includegraphics[width=0.8\textwidth]{screenshots_300dpi/game/mask_tiles_3.png}
 	\caption{\cw{OR} foreground bitmap.}
 \end{subfigure} 

\caption{Draw masked foreground tile.}
\label{fig:foreground_tile}
\end{figure}

\par
Combining the background and foreground layers involves using a bitwise \cw{AND}-operation to clear the background pixels, followed by a bitwise \cw{OR}-operation to write the foreground pixels. Because each tile is 2 bytes wide, writes to VRAM therefore always occur on even memory addresses, ensuring full utilization of the 80286 CPU's 16-bit data bus.\\

\par
\begin{minipage}{\textwidth}
  \lstinputlisting[language={[x86masm]Assembler}]{code/update_tile.asm}
  \end{minipage}
  \label{update_tile}


\subsection{Drawing sprites}
\label{section:drawing_sprites}
The next step is to render sprites on the screen. Most home computers of that era had built-in sprite functionality on the video card. For example, on an MSX computer, one could simply enter\\

\par
\begin{minipage}{\textwidth}
  \lstinputlisting[language=C]{code/msx_sprite.c}
  \end{minipage}
  \label{msx_sprite}

\par
to display a sprite on the screen. Updating the \cw{(X, Y)} coordinates would move the sprite, with the display adapter handling everything else. Unfortunately, the concept of sprites did not exist on EGA cards, so game developers had to implement their own solution.\\

\par
A challenge arises from the fact that sprites can move freely across the screen and are not byte-aligned. To address this, the bit-shifting technique described in section "\nameref{section:bitshifting}" on page \pageref{section:bitshifting} is used. When caching a sprite in memory, each sprite is stored four times, with each copy shifted by two or more pixels. The property \cw{*spr->shifts} determines the number of bit shifts applied to each of the four copies.\\


\begin{figure}[H]
  \centering
  \includegraphics[width=\textwidth]{imgs/drawings/sprite_shift.eps}
  \caption{Sprite shifted in 4 steps.}
  \label{fig:sprite_shift}  
\end{figure}

\par
Displaying the correct shifted sprite is as simple as\\
\par
\begin{minipage}{\textwidth}
  \lstinputlisting[language=C]{code/sprite_shift.c}
  \end{minipage}
  \label{sprite_shift}  
  
\par
Creating and storing the bit-shifted sprites takes place during loading and caching the sprites from disk. The first sprite keeps the original width (\cw{smallplane}), every shifted copy is one byte wider to accommodate the 8-bit shift (\cw{bigplane}). 
\\
  
\par
\begin{minipage}{\textwidth}
  \lstinputlisting[language=C]{code/shift_sprite.c}
  \end{minipage}
  \label{sprite_shift}  

\par
One might ask why only four shifts are used instead of eight to achieve a perfect one-pixel shift. Consider Commander Keen, which uses 91 different sprites that are all loaded permanently into memory. These sprites consume 41,095 bytes of RAM, and after creating seven additional shifted copies, this grows to a total of 328,760 bytes. In addition, the engine still needs to load and copy all other sprites required for a level. This memory cost was deemed too high. To save RAM, the number of bit-shifted copies is therefore limited to four, resulting in a two-pixel alignment.\\

\par
Recall that the map itself scrolls with one-pixel alignment using the EGA Pel Pan register. To keep the sprites and the world synchronized, pixel panning must therefore also be restricted to two pixels. A simple technique is used to enforce this constraint: the least significant bit of the Pel Pan value is cleared using a mask.\\

\par
\begin{minipage}{\textwidth}
  \lstinputlisting[language=C]{code/pixel_mask.c}
  \end{minipage}
  
  
\par
If multiple sprites are displayed on the same tile, each sprite is assigned a priority from 0 to 3 to determine the drawing order. A sprite with a higher priority number is always drawn on top of sprites with a lower priority. Since sprites are always drawn on top of tiles, this can create unnatural situations, such as when Commander Keen is climbing through a hole, as illustrated in Figure \ref{fig:draw_layers}.\\

\setlength{\fboxsep}{0pt}
\begin{figure}[H]
\begin{subfigure}{.25\textwidth}
  \centering
  \fbox{\includegraphics[width=.9\textwidth]{screenshots_300dpi/game/tile_composite_1.png}}
  \caption{Background tile.}
\end{subfigure}%
\begin{subfigure}{.25\textwidth}
  \centering
  \fbox{\includegraphics[width=.9\textwidth]{screenshots_300dpi/game/tile_composite_2.png}}
  \caption{Foreground tile.}
\end{subfigure}
\begin{subfigure}{.25\textwidth}
  \centering
  \fbox{\includegraphics[width=.9\textwidth]{screenshots_300dpi/game/tile_composite_3.png}}
  \caption{Sprite on top.}
\end{subfigure}
\caption{Unnatural situation where Commander Keen is in front of a hole.}
\label{fig:draw_layers}
\end{figure}

\par
To draw sprites 'inside' a foreground tile, a trick is used by introducing a priority foreground tile in the \cw{tinf} info table. The attribute is named INTILE. If a foreground tile has its INTILE attribute high bit set (\cw{80h}), sprites with priority 0--2 will not be drawn over it. This results in the following drawing order:
\begin{enumerate}
  \item Draw the background tile and the masked foreground tile.
  \item Draw sprites with priority 0, 1, and 2 (in that order), and mark the corresponding tiles in the tile buffer array with a '3', as illustrated in Figure \ref{fig:kc4_6_step4} on page \pageref{fig:kc4_6_step4}. 
  \item Scan the tile buffer array for tiles marked with '3'. If the corresponding foreground tile's INTILE attribute high bit (\cw{80h}) is set, re-draw the foreground tile. 
  \item Finally, draw sprites with priority 3. These sprites are always drawn on top of everything.
\end{enumerate}
\par

\setlength{\fboxsep}{0pt}
\begin{figure}[H]
\centering
\begin{subfigure}[t]{.245\textwidth}
  \centering
  \fbox{\includegraphics[width=.9\textwidth]{screenshots_300dpi/game/tile_composite_1.png}}
  \caption{Background tile.}
\end{subfigure}%
\begin{subfigure}[t]{.245\textwidth}
  \centering
  \fbox{\includegraphics[width=.9\textwidth]{screenshots_300dpi/game/tile_composite_2.png}}
  \caption{Foreground tile.}
\end{subfigure}
\begin{subfigure}[t]{.245\textwidth}
  \centering
  \fbox{\includegraphics[width=.9\textwidth]{screenshots_300dpi/game/tile_composite_3.png}}
  \caption{Sprite on top.}
\end{subfigure}
\begin{subfigure}[t]{.245\textwidth}
  \centering
  \fbox{\includegraphics[width=.9\textwidth]{screenshots_300dpi/game/tile_composite_4.png}}
  \caption{Redraw foreground tile.}
\end{subfigure}
\caption{Draw sprite inside a tile, by redrawing foreground tile.}
\label{fig:clip_tinf}
\end{figure}

\par
\begin{minipage}{\textwidth}
  \lstinputlisting[language={[x86masm]Assembler}]{code/mask_foreground_tile.asm}
\end{minipage}
\label{mask_foreground_tile}
\par

\subsection{Tile Draw Performance Tricks}
The minimum number of read/write operations per tile, when only a background tile is required, is 128 bytes (2$\times$16 bytes$\times$4 memory banks). Worst case, up to 512 bytes of read/write operations are required (background tile, foreground tile, sprite, and foreground tile again), with multiple bitwise operations involved.
The engine employs several tricks to squeeze the maximum out of each CPU cycle: background tile caching and word-aligned memory writing.

\subsubsection{Background Tile Caching}
When updating the master screen in VRAM, the engine copies each tile from RAM to VRAM. Copying each pixel involves one read and one write operation across the four memory banks. In section "\nameref{section:optimize_tile}" on page \pageref{section:optimize_tile} it was explained how reprogramming EGA latches enables copying four bytes at once between VRAM locations, a technique that can be applied as well to tiles containing only a background layer.\\

\par
Once a background tile is loaded into the master screen, subsequent requests for the same tile can be handled by copying it directly from VRAM using the reprogrammed latches, rather than copying from RAM. The engine only needs to track which background tiles are already loaded into VRAM using an array.\\

\par 
\begin{minipage}{\textwidth}
  \lstinputlisting[language=C]{code/tile_cache.c}
\end{minipage}

\par
The assembly function \cw{RFL\_NewTile} is responsible for drawing tiles to the master screen.\\

\par
\begin{minipage}{\textwidth}
  \lstinputlisting[language={[x86masm]Assembler}]{code/RFL_NewTile.ASM}
\end{minipage}
\par

\par
It first checks whether the background tile is already stored in the tile cache array. If it is, the tile is copied at 32 bits per cycle from VRAM to VRAM. Otherwise, the tile is loaded from RAM, and the VRAM pointer for that tile is stored in the tile cache array for future use.


\subsubsection{Writing to word-aligned memory}
The 286 CPU can read and write 16 bits in a single cycle, but there is a caveat: writing a word at offset \cw{FFFFh} causes the CPU to crash.
This is not an issue for tiles, because they are word-aligned (16 bits wide). Since the screen scrolls in tile-sized steps, they always reside at even memory addresses. \\

\par
Sprites, however, are byte-aligned. When a sprite is drawn starting at an odd memory address, it may cross the \cw{FFFFh} offset and crash the CPU. To avoid this, the engine uses small, specialized routines for each combination of sprite width and even/odd alignment, ensuring correct wrapping at \cw{FFFFh}. \\

\par
For example, when drawing a 6-byte-wide sprite that starts at an even address, three \cw{word} writes wrap around the segment. If it starts at an odd address, the routine writes a \cw{byte}, followed by two \cw{word} writes, and finally another \cw{byte}. \\

\begin{figure}[H]
  \centering
  \includegraphics[width=0.8\textwidth]{imgs/drawings/mask_block.eps}
  \caption{Word-optimized 6 byte wide sprite to even (\cw{MASK6E}) and odd (\cw{MASK6O}) address.}
  \label{fig:mask_block}  
\end{figure}

\par
In total, 23 routines handle aligned \cw{word}-sized writes for sprites up to 10 bytes wide. Their function pointers are stored in an array. \\


\par
\begin{minipage}{\textwidth}
  \lstinputlisting[language={[x86masm]Assembler}]{code/VW_maskroutines.asm}
\end{minipage}

\par
The engine first determines whether the destination is at an even or odd memory address. If the address is odd, it first writes a single byte to align the subsequent operations to an even address, after which it continues writing in 16-bit words. 
By cleverly using bit-shift operations and the Carry Flag (CF), the engine calls the appropriate function to write the sprite to VRAM.\\

\par
\begin{minipage}{\textwidth}
  \lstinputlisting[language={[x86masm]Assembler}]{code/VW_MaskBlock.asm}
\end{minipage}
\pagebreak





\end{document}
